\documentclass[11pt,a4paper]{article}
\usepackage[T1]{fontenc}
\usepackage{lmodern,amsmath,amssymb,amsthm,mathtools,booktabs,microtype}
\usepackage[margin=27mm]{geometry}
\usepackage[colorlinks=true,linkcolor=blue,citecolor=blue,urlcolor=blue]{hyperref}
\setlength{\emergencystretch}{2em}
\newtheorem{theorem}{Theorem}[section]
\newtheorem{lemma}[theorem]{Lemma}
\newcommand{\Q}{\mathbb Q}\newcommand{\Z}{\mathbb Z}
\DeclareMathOperator{\Tr}{Tr}
\title{An inert obstruction to a deterministic contact coupling}
\author{}\date{17 November 2026}
\begin{document}\maketitle
\begin{abstract}
An order-three symmetry of a six-contact function identifies two
contact square classes and suggests eliminating their common free
slope. We prove that a specified deterministic height substitution
fails for every positive odd coefficient shift. At one fixed inert
prime, infinitely many positive integer inputs have simple contacts
in both blocks. Every guarded auxiliary tuple then has a forbidden
negative odd detector pole. A finite certificate of thirteen rows
and a formal-group lifting argument establish the uniform assertion.
The obstruction is confined to the displayed substitution with its
prescribed root index.
\end{abstract}

\section{The coupling and the obstruction}

Set
\[
 K=\Q(\sqrt2),\quad D=11-6\sqrt2,\quad E:y^2=x^3-Dx,
 \quad P=\left(\frac{57-28\sqrt2}{8},
                 \frac{-660+417\sqrt2}{32}\right).
\]
For positive integers \(X\) and a fixed positive odd integer \(B\), put
\[
 A=[X]P,\quad T_X=\Tr_{K/\Q}x(A),\quad
 u=(X^2+B)^4+T_X^4,
\]
\begin{equation}\label{eq:coupling}
 t_0=\frac{1-3u}{7u},\qquad t_1=-\frac{u+2}{7},\qquad
 a=\frac{u(u+2)}{3u-1}.
\end{equation}
Define
\[
 F(t)=\frac{t(t+1)(2t+1)(3t+1)(4t+1)(5t+2)}{120},
 \quad q(t)=7t^2+5t+1,\quad \eta(t)=\frac{F(t)}{q(t)^3}.
\]
The point and difference tests are those of the conditional source
\cite{Conditional,Reweighted}. In particular, the detector product
contains \(f(H)\) and \(\overline{f(H-P)}\), where \(f=1/x\).
All denominators and the reciprocal trace are guarded.

\begin{theorem}\label{thm:main}
For every positive odd \(B\), infinitely many positive integers \(X\)
satisfy
\[
 v_{81421}(F(t_0))=v_{81421}(F(t_1))=1.
\]
Every displayed denominator and scale guard is a unit at this prime.
For either choice
\[
 (d_0,w)=\left(F(t_0),\frac{F(t_0)}{4(t_0-1)^8}\right)
 \quad\text{or}\quad
 (d_0,w)=\left(\eta(t_0),
                  \frac{\eta(t_0)}{(1+\eta(t_0))^2}\right),
\]
every defined auxiliary tuple in
\[
 W=\prod_i(1+C_i^2/d_0),\qquad V=\Tr(W^{-1}),\qquad b=4w/V^2
\]
has \(v_{81421}(b)<0\) odd, or violates the guard \(V\ne0\).
Consequently this prescribed-index deterministic coupling cannot give
uniform completeness for any pole test excluding such valuations.
\end{theorem}

The prime is outside the conditional exceptional set. The conclusion
allows arbitrary auxiliary points; it is not limited to multiplication
families. It does not apply to a different height substitution,
an independently varying root label, or a construction retaining a
free slope.

\section{Contact identities and the inert-place lemma}

The transformation \(\gamma(t)=-(2t+1)/(7t+3)\) has order three and
\[
 F(\gamma(t))=\frac{F(t)}{(7t+3)^6},\qquad
 q(\gamma(t))=\frac{q(t)}{(7t+3)^2}.
\]
Thus \(\eta\circ\gamma=\eta\) and \(t_1=\gamma(t_0)=at_0\).
Substitution gives
\begin{align*}
 F(t_0)&=\frac{(u-5)(u+2)(2u-3)(3u-1)(4u+1)(5u-4)}
                    {14117880u^6},\\
 F(t_1)&=u^6F(t_0),\qquad
 q(t_0)=\frac{u^2-u+1}{7u^2},\qquad q(t_1)=\frac{u^2-u+1}{7}.
\end{align*}
For \(u>5\), both contact values are positive and the real domain
guards are nonzero.

\begin{lemma}\label{lem:inert}
Let \(p\ne2,7\) be inert in \(K\), and suppose
\(v_p(d_0)=v_p(w)=1\). A defined product containing the two tests
\(f(H),\overline{f(H-P)}\) has \(V\in p\Z_p\).
When \(V\ne0\), the valuation of \(4w/V^2\) is negative and odd.
\end{lemma}
\begin{proof}
The curve has good reduction, since its discriminant is \(64D^3\)
and \(\operatorname{N}_{K/\Q}(D)=49\). The affine coordinates of
\(P\) are integral and its reduction is not the identity.
Conjugation preserves the unique valuation above an inert prime.
If both displayed tests had positive valuation, both \(H\) and
\(H-P\) would reduce to the identity, forcing \(P\) to do so.
Hence at least one test has nonpositive valuation.

Every factor is nonzero and has valuation
\[
 v(1+C_i^2/d_0)=\min(0,2v(C_i)-1)\le0;
\]
the two terms have different valuation parity. A zero test gives
factor one. At least one factor has valuation at most \(-1\),
so \(W^{-1}\in p\mathcal O_{K_p}\) and \(V\in p\Z_p\).
The last assertion follows from \(v_p(4w/V^2)=1-2v_p(V)\).
\end{proof}

\section{The finite certificate}

Put \(p=81421\), \(N=6629217184\).
Exact arithmetic verifies that \(p\) is prime,
\(p\equiv5\pmod8\), \(p\nmid N\), and
\([N]\widetilde P=O\) over \(\mathbb F_{p^2}\).
Only the annihilator property, not an exact order, is used.
The following rows satisfy
\[
 T_n=\Tr x([n]\widetilde P),\qquad y^4+T_n^4=u_0
       \quad\text{in }\mathbb F_p.
\]
\begin{center}
\begin{tabular}{rrrr}\toprule
\(n\)&\(T_n\)&\(u_0\)&\(y\)\\\midrule
1&61080&\(4/5\)&8576\\
1&61080&\(4/5\)&35004\\
6&17254&\(4/5\)&73321\\
1&61080&\(4/5\)&46417\\
10&74125&\(3/2\)&3133\\
9&45107&\(4/5\)&80711\\
20&31595&5&57446\\
19&25083&5&42108\\
23&78190&\(3/2\)&73\\
18&57250&\(3/2\)&1060\\
29&41438&\(3/2\)&77025\\
24&26788&\(3/2\)&23167\\
2&21974&5&29234\\\bottomrule
\end{tabular}
\end{center}
Writing \(Y\) for the last column, enumeration of the finite field gives
\begin{equation}\label{eq:coverage}
 \bigcup_{y\in Y}\bigl(y-(\mathbb F_p^\times)^2\bigr)=\mathbb F_p.
\end{equation}
All listed ordinates, abscissas and \(y\)'s are units. Each of the
three values \(u_0\) is a simple root of \(F((1-3u)/(7u))\).
All slope, scale and denominator guards are units at these roots.
The distributed Python and independent integer-arithmetic JavaScript
checkers reconstruct the group calculations, all \(81421\) covered
residues, the simple contacts and the fixture in Section~\ref{sec:fixture}.

\section{Lifting for every coefficient shift}

\begin{proof}[Proof of Theorem~\ref{thm:main}]
Fix \(B\). By \eqref{eq:coverage}, choose a row and \(x\ne0\) with
\(x^2+B=y\pmod p\). The Chinese remainder theorem gives \(X_0>0\) with
\[
 X_0\equiv n\pmod N,\qquad X_0\equiv x\pmod p.
\]
The annihilator identity implies
\([X_0]\widetilde P=[n]\widetilde P\), so \(u\equiv u_0\pmod p\).

The point \([N]P\) is in the first formal-group filtration.
The integral multiplication series \([p](z)=pz+O(z^2)\) puts
\([pN]P\) in the second filtration. Adding \(pN\) to the index
therefore preserves the affine point coordinates and \(T_X\) modulo
\(p^2\). It changes the other summand of \(u\) by
\[
 ((X+pN)^2+B)^4-(X^2+B)^4
       \equiv8X(X^2+B)^3pN\pmod{p^2}.
\]
After division by \(p\), this has nonzero residue, since
\(x,y,N\) are units. At the simple contact, at least one of
\(X_0,X_0+pN\) therefore gives \(v_p(F(t_0))=1\).
The identity \(F(t_1)=u^6F(t_0)\), with \(u\) a unit, gives
the second valuation.

Adding positive multiples of \(p^2N\) preserves these congruences
and gives infinitely many positive \(X\), eventually with \(u>5\).
All the denominator-unit statements persist. In particular,
\(v_p(d_0)=v_p(w)=1\) for each specified detector choice.
Lemma~\ref{lem:inert} completes the proof.
\end{proof}

\section{An explicit input}\label{sec:fixture}

For \(X=98,B=81\), calculation modulo \(p^2=6629379241\) gives
\[
 x(A)=1181463837+1902694508\sqrt2,\qquad
 y(A)=3842721784+4607123365\sqrt2,
\]
\[
 \begin{gathered}
 T_X=2362927674,\quad u=4429750217,\\
 t_0=4664473388,\quad t_1=2208341072,\quad a=4421038170.
 \end{gathered}
\]
In particular,
\[
 F(t_0)\equiv44771p,\qquad F(t_1)\equiv73604p\pmod{p^2}.
\]
The same input works for all odd shifts
\(B=81+13258758482k\), \(k\ge0\). The coverage and lifting
argument proves the stronger assertion for every positive odd shift.

\begin{thebibliography}{9}\raggedright
\bibitem{Conditional} \emph{Thirteen Variables from Direct Elliptic
Index Comparison}, repository manuscript,
\nolinkurl{papers/conditional_few_variables/direct_input13_20261110_v22.tex}.
\bibitem{Reweighted} \emph{Quadratic reweightings of native elliptic detectors}, repository
manuscript, 16 November 2026,
\nolinkurl{papers/pole_parameters/reweighted_elliptic_detectors_20261116_v53.tex}.
\end{thebibliography}
\end{document}
